\section{The answer key: a training-objective OCP}
\label{sec:answer}
% Code: src/apexzero/ocp/window.py (TO-OCP, window mode); src/apexzero/ocp/replay.py and
% scripts/spike_to_ocp.py (replay check). Plan: docs/plan/two-track-plan.md sections 1, 2.2.
% Evidence: docs/findings/to-ocp-spike.md, docs/STATUS.md (K2a row).

\subsection{Problem}
\label{sec:answer-problem}

For a level $\theta$ with start state $x_0$ and window length $T = 1200$ control steps,
the training-objective OCP is
\begin{equation}
\begin{aligned}
  \Jstar(\theta) \;=\; \max_{a_0,\dots,a_{T-1}} \quad
    & \sum_{t=0}^{T-1} \gamma^{t}\Big[\frac{s_{t+1}-s_t}{10\,\text{m}}
      - 0.02\,\lVert a_t - a_{t-1}\rVert^2\Big] \\
  \text{s.t.}\quad
    & x_{t+1} = \Phi_\theta(x_t, a_t), \qquad a_t \in [-1,1]^2, \qquad a_{-1} = 0, \\
    & -\big(w_r(s_{t+1}) - 1\,\text{m} - \epsilon\big) \le n_{t+1}
      \le w_l(s_{t+1}) - 1\,\text{m} - \epsilon, \\
    & |\alpha_f(x_{t+1})|,\ |\alpha_r(x_{t+1})| \le \alpha_{\max}, \qquad
      |\xi_{t+1}| \le \pi/2, \qquad v_{x,t+1} \ge 2\,\text{m/s},
\end{aligned}
\label{eq:toocp}
\end{equation}
for $t = 0, \dots, T-1$, with $\gamma = 0.995$.
% src: src/apexzero/ocp/window.py (module docstring; WindowSpec: steps 1200, gamma 0.995,
%      edge_margin 1.0, safety 0.03; WindowSolver bounds)
Here $\Phi_\theta$ is the simulator's own control-step map: RK2 with ten 5\,ms substeps
and the action held for 50\,ms, the same vehicle ODE, and the level's six-feature car
context as a parameter. The solver has no separate grip input, so the hidden grip must
be folded into the context: the solver receives the context with its global friction
scale multiplied by $\mu_h$, which is how the environment builds its parameters from
$(c, \mu_h)$.
% src: src/apexzero/ocp/window.py (_context_params: ctx[0..5] only; WindowSolver.solve ctx);
%      src/apexzero/env/racing.py (params_from_context: ctx[4] * hidden_mu)
% src: src/apexzero/ocp/window.py (WindowSolver.solve(hidden_mu=): ctx_drive[4] *= hidden_mu)
\TODO{add a test that the solver's step map with folded grip equals the environment's step
for $\mu_h \ne 1$}
The constraints are the environment's non-termination set (\cref{sec:env-mdp}), applied
at every post-step state and tightened in three ways that can only lower the optimum.
First, a safety margin $\epsilon = 3$\,cm keeps a float64 replay of a solution that rides
a constraint from terminating. Second, the environment terminates only after the car has
stayed below 2\,m/s for 2\,s, while the OCP requires $v_x \ge 2$\,m/s at every step.
Third, the solver bounds $v_x \le 120$\,m/s, $|v_y| \le 25$\,m/s and $|r| \le 5$\,rad/s.
The objective is the environment's reward \eqref{eq:reward}; the crash penalty does not
appear because terminating trajectories are infeasible.
% src: src/apexzero/ocp/window.py (WindowSolver.solve: lbX/ubX, v_x lower bound 2.0);
%      src/apexzero/env/racing.py (step_no_reset: slow-for-2-s termination)
\TODO{log whether the stall bound and the auxiliary state bounds are ever active in the
label runs}

\paragraph{Exact track lookup} CasADi's B-spline node exists only for its MX graph type,
which made exact Hessians about 50 times slower. Instead, each control step $t$ carries a
window of $K$ spline coefficients around its warm-start position $\sigma_t$, and
$\kappa(s)$ and $w(s)$ are evaluated as the exact B-spline sum over that window.
Inside the window this equals the simulator's spline; $s_t$ is bounded to within 15\,m of
$\sigma_t$, and an outer loop re-centres the windows and re-solves if a solution reaches
the bound. Track coefficients, car context and $x_0$ are solver parameters, so one
compiled solver serves every level with the same window length.
% src: src/apexzero/ocp/window.py (module docstring; WindowSpec.radius = 15.0; solve loop)
The CasADi step map and windowed lookup agree with the JAX simulator to
$5\times10^{-17}$ in the state and $2\times10^{-17}$ in the lookup on a development
circuit (the test asserts $10^{-10}$ and $10^{-12}$).
% src: tests/test_window_ocp.py; paper/tools/gate_values.py (Hockenheim)
The NLP is a multiple-shooting problem with every state and action as a decision
variable ($8(T+1) + 2T$ of them) and one defect constraint per control step, solved by
IPOPT with the MUMPS linear solver~\citep{wachter2006ipopt,andersson2019casadi}.
% src: src/apexzero/ocp/window.py (WindowSolver.__init__)

\paragraph{Warm starts and labels} Every level gets the same three starts: the $s$-domain
minimum-time optimum mapped to time, the quasi-steady-state line, and the best rollout of
the final checkpoint. The label is the best value over the starts whose replay passes the
check below. The budget is uniform because a second pass only on suspicious levels would
bias where labelling effort goes.
% src: docs/plan/two-track-plan.md sections 2.2 ("Every level gets the same 3-start
%      budget") and 9 (scope critic: uniform agent-warm-start pass)

\paragraph{Replay verification} The optimal controls are applied open-loop, one action per
50\,ms, in the float64 JAX environment with its termination checks. A label is accepted
only if the replay survives the whole window and its discounted return matches the
solver's value to within 0.05\%.
% src: docs/plan/two-track-plan.md section 2.2 ("Replay in float64 must match within
%      0.05%, with a 2-5 cm margin inside the corridor"); docs/findings/to-ocp-spike.md
This closes the loop between the two compilations: if the CasADi model and the simulator
disagreed anywhere along the trajectory, the replay would drift or terminate.

\subsection{Why this is the right answer key}
\label{sec:answer-why}

UED's minimax-regret objective refers to $\max_{\pi'} J^{\pi'}(\theta) - \Jpi(\theta)$
with $J$ the agent's own expected discounted return. \Cref{eq:toocp} is that maximum for
the agent's problem, with four properties that a generic reference (a lap-time optimum,
a success indicator, a point-mass model) lacks.
\begin{enumerate}
  \item \textbf{Same objective.} Regret comes out in the agent's units, discounted
        return, which is the quantity the level scores are meant to track. A time-based
        reference is not in those units: turning a lap-time or distance gap into a return
        gap needs a conversion that depends on where in the window the time is lost, and
        that conversion would set the scale of every regret we report. The discount
        weighting the agent sees is the one the optimum sees. In one dimension
        (\cref{app:braking}) the two optima brake at the same point, to within one control
        step, at the agent's $\gamma = 0.995$ and diverge only at horizons of a few tenths
        of a second; the appendix is mainly an independent check that a discounted optimum
        built from the environment's reward and crash penalty behaves as expected.
  \item \textbf{Same discretisation and control rate.} The optimiser cannot exploit a finer
        integrator or a faster controller than the agent has, so integration error and
        the cost of acting at 20\,Hz do not appear as regret.
  \item \textbf{Same actuators, and constraints that are never looser.} The steering and
        pedal lags, the rate limit, the corridor and the slip limit are all inside
        $\Phi_\theta$ or the constraints, so the optimum is reachable by a policy acting
        through the same actuators. Where the constraints differ from the environment's
        termination rules (the margin, the stall bound, the auxiliary bounds) they are
        stricter.
  \item \textbf{Determinism per level.} Given $\theta$, the dynamics are deterministic
        (the hidden grip is fixed per level), so the best open-loop control sequence from
        $x_0$ achieves the optimal value over all policies, closed-loop or not. Optimising
        over open-loop sequences loses nothing.
\end{enumerate}
The fourth point also says what $\Jstar$ is \emph{not}: it is the optimum for an agent that
knows $\theta$. Our agent sees a 150--500\,m preview, noisy observations and a car
context, but not the hidden grip. The part of regret caused by what the policy cannot
observe is irreducible for this policy class. That is regret in the sense UED uses it,
but it is exactly the kind a regret-maximising sampler can fixate
on~\citep{beukman2024remidi}; we return to it in \cref{sec:threats}.

\subsection{A lower bound, not an estimate}
\label{sec:answer-lb}

Let $J^{\mathrm{opt}}(\theta)$ be the supremum of the window return over all control
sequences, including those that terminate. Three things can make the label fall short of
it: IPOPT returns a local optimum; the feasible set is tightened by the safety margin, the
stall bound and the auxiliary state bounds; and terminating trajectories are excluded,
although with
$\gamma^{1200} \approx 0.0024$ a late crash can in principle cost less than the progress it
buys.
% src: gamma^1200 ~ 0.0024 as stated in docs/paper/sources-jaxued.md section 4 item 3
Nothing can make it exceed $J^{\mathrm{opt}}$ if the label is the \emph{replayed} return,
because a replay is a trajectory the simulator actually produced; the solver's own value
agrees with it to within the 0.05\% acceptance tolerance.
\TODO{fix in \code{audit-v1} that the stored label is the float64 replayed return, not the
solver's objective value}
Hence $\Jstar \le J^{\mathrm{opt}}$, and for every policy
\begin{equation}
  \Reg(\theta) \;:=\; \Jstar(\theta) - \Jpi(\theta) \;\le\;
  J^{\mathrm{opt}}(\theta) - \Jpi(\theta).
  \label{eq:lowerbound}
\end{equation}
We therefore call $\Jstar$ the \emph{best-known optimum} and $\Reg$ a \emph{lower bound}
on regret throughout. Any return the agent itself realises is also a lower bound on
$J^{\mathrm{opt}}$ (up to the rounding of the single-precision rollouts), because a
realised action sequence is a feasible open-loop control; the best observed agent return
is thus a free check on the label. We report how often it
exceeds the label: \RESULT{fraction of levels, per phase, on which an observed agent
return beats $\Jstar$, and by how much}. The bound holds for the expectation $\Jpi$; a
Monte Carlo estimate $\Jhat$ can still put individual levels below zero.

Two references bracket the label from other sides. A control-rate reference $\Jhund$
solves the same window at a 100\,Hz zero-order hold, warm-started from the 20\,Hz
solution, on 200 levels; $\Jhund - \Jstar$ is the price of acting at 20\,Hz, which no
20\,Hz policy can recover. An undiscounted ($\gamma = 1$) maximum-distance variant on 600
A\textsubscript{DR} levels gives distance lost in 60\,s without discount weighting.
% src: docs/plan/two-track-plan.md section 2.2 ("Bounds and references", "Secondary targets")
Label noise is measured directly: 200 pilot levels get at least three extra starts, which
gives a per-level label standard deviation (\cref{sec:protocol-targets}).

\paragraph{The window objective and the agent's training objective} PPO bootstraps the
value at the window end, so the agent is trained for the infinite-horizon discounted
return, while \cref{eq:toocp} stops at $T$. The omitted tail is $\gamma^{T}V(x_T)$, with
$\gamma^{1200} \approx 0.0024$. The same finite horizon lets the OCP spend the last
seconds of a window in ways a time-unaware policy cannot (for example by not slowing for a
corner just beyond the window end); that gain is discounted by roughly the same factor.
We expect both effects to be small and quantify them in \cref{sec:threats}.

\subsection{Evidence so far}
\label{sec:answer-evidence}

The spike solved the TO-OCP on a 10\,s window of a development circuit (Hockenheim),
warm-started from the $s$-domain minimum-time optimum mapped to time. IPOPT converged
(\code{Solve\_Succeeded}) in 169 iterations over 3 re-centring rounds, in 44\,s of solve
time and 0.7\,GB of peak memory in a single process on a 12-thread laptop CPU. The
float64 replay reproduced the
optimum exactly ($J^* = 48.089$ both ways) and survived the window.
% src: docs/findings/to-ocp-spike.md (table)
Two very different warm starts reached the same optimum.
% src: docs/STATUS.md (K2a row)
These are one window of one circuit, a sixth of the label window length. Whether the
solver converges as reliably, and how long it takes, on 60\,s windows of procedural
strips decides the size of the label budget.
On eight production-like generated levels (8\,km strips, each with its own random car and
hidden grip), a 60\,s window, the level's own $s$-domain optimum included, took 131--520
CPU-seconds (median 206) on seven; the eighth hit IPOPT's iteration limit after 791
CPU-seconds with no verified optimum. Windows of Hockenheim took 106--189 CPU-seconds, and
every verified optimum replayed to its $J^*$ within $3\times10^{-12}$.
% src: docs/findings/to-ocp-spike.md (production-like levels pool:0-7; Hockenheim windows);
%      the benchmark's other circuit windows lie on lockbox circuits and are sealed
\RESULT{convergence rate and per-solve CPU time over the production label set}
The plan budgets 1,200--3,000 CPU-hours for about \num{5400} levels with three starts
each. Three starts sharing one $s$-domain optimum cost about 775 CPU-seconds per level at
these figures, about \num{1160} CPU-hours for \num{5400} levels; above 12 CPU-minutes per
solve the pool would shrink to \num{2000} levels and the curriculum buffer to 500.
% src: docs/plan/two-track-plan.md sections 5.2 (Spike gate) and 8 (CPU item)
